\chapter{Future Work}

\hspace{\parindent} In the future, the author aims to integrate within the current \textit{NeuroSEE} pipeline a novel information theoretic and machine learning-based algorithm for reducing the dimensionality of the large-scale multivariate signals extracted, as well as an algorithm that is capable of estimating information flow graphs that reflect disease progression. 

For the place field analysis, the author aims to further improve the algorithm by using information theoretic approaches in accurately estimating the spatial information content, e.g. using Gaussian Process-based method. It would also be helpful if other information theoretic analysis are integrated within the pipeline, e.g. spike inference and spike entropy calculation, in order to quantify information encoding and transmission.

On a more general note, the author aims to use the proposed \textit{NeuroSEE} pipeline in investigating how the spatiotemporal dynamics of neural cortical circuits involved in memory encoding and recall are affected by Alzheimer's disease as it progresses. More specifically, it would be important to characterise the differences between different mouse models of AD in terms of their effects on information flow in hippocampal-cortical circuits, so as to generalise our understanding on the effect of the different manifestations of AD on memory circuit dynamics. 

In mouse models of AD, amyloid or tau protein-related pathologies result in changes in cortical circuitry, which are particularly evident in the hippocampus. What is not known is how these changes relate to the function of this circuitry in memory encoding and recall processes, or how therapies which have been found to restore behavioural performance in mouse models affect this circuitry. Hence, it is good to examine the effects of therapeutic interventions, including pharmacological agents and neuromodulatory therapies, on the information processing properties in different mouse models. 

All of the above-mentioned works will be performed as the author progresses towards the PhD phase of the Centre for Doctoral Training in Neurotechnology for Life and Health programme. 